% ============================================================================
%  styles/boxes.tex — 学术环境（书籍形制）
%
%  与原模板的差别（书籍化改造的核心）：
%    · 原版六大定理盒的标题栏是「莱茵生命 logo 戳 + 竖排中文 + 英文粗体」
%      预渲染矢量图。那是游戏 IP 视觉，出版物不能带，因此这里改为
%      「粗左线 + 彩色标签 + 编号 + 可选中注」的学术书籍形制，
%      六色编码（定义 / 定理 / 命题 / 推论 / 引理 / 例题）完整保留 ——
%      这才是孤星配色系统真正有价值的部分：读者扫一眼即可分辨环境类型。
%    · 定理环境直接用 tcolorbox + 计数器实现，不再走 amsthm 环境二次包装，
%      计数器形制仍为 \thechapter.\arabic{...}，\label / cleveref 行为不变。
% ============================================================================

% amsthm 已定义 \proof；本类自定义 proof 环境，先解除占用（\let 为 \relax 后
% xparse 视为未定义，可重新 \NewDocumentEnvironment）
\let\proof\relax
\let\endproof\relax

% --- 中文环境名 -------------------------------------------------------------
\providecommand{\definitionname}{定义}
\providecommand{\theoremname}{定理}
\providecommand{\propositionname}{命题}
\providecommand{\corollaryname}{推论}
\providecommand{\lemmaname}{引理}
\providecommand{\axiomname}{公理}
\providecommand{\postulatename}{公设}
\providecommand{\propertyname}{性质}
\providecommand{\examplename}{例题}
\providecommand{\exercisename}{习题}
\providecommand{\solutionname}{解}
\providecommand{\remarkname}{注}
\providecommand{\proofname}{证明}

% --- 计数器（各环境独立，按章编号；编号形制与 elegantbook 时代一致）---------
\newcounter{ltdefinition}[chapter]
\renewcommand{\theltdefinition}{\thechapter.\arabic{ltdefinition}}
\newcounter{lttheorem}[chapter]
\renewcommand{\thelttheorem}{\thechapter.\arabic{lttheorem}}
\newcounter{ltproposition}[chapter]
\renewcommand{\theltproposition}{\thechapter.\arabic{ltproposition}}
\newcounter{ltcorollary}[chapter]
\renewcommand{\theltcorollary}{\thechapter.\arabic{ltcorollary}}
\newcounter{ltlemma}[chapter]
\renewcommand{\theltlemma}{\thechapter.\arabic{ltlemma}}
\newcounter{ltaxiom}[chapter]
\renewcommand{\theltaxiom}{\thechapter.\arabic{ltaxiom}}
\newcounter{ltproperty}[chapter]
\renewcommand{\theltproperty}{\thechapter.\arabic{ltproperty}}
\newcounter{ltexample}[chapter]
\renewcommand{\theltexample}{\thechapter.\arabic{ltexample}}
% 习题全章共用一个计数器（exercise / exercises / Exercises 连续编号）
\newcounter{bookexercise}[chapter]
\renewcommand{\thebookexercise}{\thechapter.\arabic{bookexercise}}
\renewcommand{\theHbookexercise}{Exercise.\thechapter.\arabic{bookexercise}}
% 注与题解收集
\newcounter{ltremark}[chapter]
\renewcommand{\theltremark}{\thechapter.\arabic{ltremark}}
\newcounter{exsolcount}[chapter]

% --- 盒子基样式 -------------------------------------------------------------
% 印刷安全：底色为该定理色的极淡 tint（≤8%），不是大面积色块。
\tcbset{
  ltbox/.style={
    enhanced,
    sharp corners,
    breakable,
    boxrule=0pt,
    leftrule=3pt,
    rightrule=0pt,
    toprule=0pt,
    bottomrule=0pt,
    top=7pt,
    bottom=7pt,
    left=10pt,
    right=10pt,
    before skip=11pt,
    after skip=11pt,
    fontupper=\normalfont,
    halign upper=justify
  },
  % 例题 / 解答 / 习题这类「内容盒」底色更淡，右侧留呼吸
  ltcontent/.style={
    ltbox,
    leftrule=2.2pt,
    top=6pt,
    bottom=6pt,
    before skip=9pt,
    after skip=9pt
  },
  % 长注 / remark* 的静音盒（键名集中在这里定义，expl3 语境里空格会被吃掉，
  % 不能直接在 \cs_new_protected 的体内写 `before skip=9pt` 这类带空格的键）
  ltremarkbox/.style={
    ltbox,
    colframe=InkMuted,
    colback=InkWash!55!white,
    leftrule=5pt,
    top=6pt,
    bottom=6pt,
    before skip=9pt,
    after skip=9pt
  }
}

% --- 定理族环境 -------------------------------------------------------------
% 各环境形制一致：粗左线（该定理色）+ 极淡 tint 底 + 彩色标签行。
% 编号用 \refstepcounter，\label / cleveref 形制与 elegantbook 时代一致。
\NewDocumentEnvironment{definition}{o}{%
  \refstepcounter{ltdefinition}%
  \LT@FloatLocal
  \begin{tcolorbox}[ltbox, colframe=definition, colback=definition!7!white]
    \noindent{\bfseries\color{definition}\definitionname~\theltdefinition
      \IfValueT{#1}{\ （#1）}}\par\nobreak\smallskip
    \ignorespaces
}{%
  \end{tcolorbox}%
}

\NewDocumentEnvironment{theorem}{o}{%
  \refstepcounter{lttheorem}%
  \LT@FloatLocal
  \begin{tcolorbox}[ltbox, colframe=theorem, colback=theorem!7!white]
    \noindent{\bfseries\color{theorem}\theoremname~\thelttheorem
      \IfValueT{#1}{\ （#1）}}\par\nobreak\smallskip
    \ignorespaces
}{%
  \end{tcolorbox}%
}

\NewDocumentEnvironment{proposition}{o}{%
  \refstepcounter{ltproposition}%
  \LT@FloatLocal
  \begin{tcolorbox}[ltbox, colframe=proposition, colback=proposition!7!white]
    \noindent{\bfseries\color{proposition}\propositionname~\theltproposition
      \IfValueT{#1}{\ （#1）}}\par\nobreak\smallskip
    \ignorespaces
}{%
  \end{tcolorbox}%
}

\NewDocumentEnvironment{corollary}{o}{%
  \refstepcounter{ltcorollary}%
  \LT@FloatLocal
  \begin{tcolorbox}[ltbox, colframe=corollary, colback=corollary!7!white]
    \noindent{\bfseries\color{corollary}\corollaryname~\theltcorollary
      \IfValueT{#1}{\ （#1）}}\par\nobreak\smallskip
    \ignorespaces
}{%
  \end{tcolorbox}%
}

\NewDocumentEnvironment{lemma}{o}{%
  \refstepcounter{ltlemma}%
  \LT@FloatLocal
  \begin{tcolorbox}[ltbox, colframe=lemma, colback=lemma!7!white]
    \noindent{\bfseries\color{lemma}\lemmaname~\theltlemma
      \IfValueT{#1}{\ （#1）}}\par\nobreak\smallskip
    \ignorespaces
}{%
  \end{tcolorbox}%
}

\NewDocumentEnvironment{axiom}{o}{%
  \refstepcounter{ltaxiom}%
  \LT@FloatLocal
  \begin{tcolorbox}[ltbox, colframe=qiancheng, colback=qiancheng!7!white]
    \noindent{\bfseries\color{qiancheng}\axiomname~\theltaxiom
      \IfValueT{#1}{\ （#1）}}\par\nobreak\smallskip
    \ignorespaces
}{%
  \end{tcolorbox}%
}

\NewDocumentEnvironment{postulate}{o}{%
  \refstepcounter{ltaxiom}%
  \LT@FloatLocal
  \begin{tcolorbox}[ltbox, colframe=qiancheng, colback=qiancheng!7!white]
    \noindent{\bfseries\color{qiancheng}\postulatename~\theltaxiom
      \IfValueT{#1}{\ （#1）}}\par\nobreak\smallskip
    \ignorespaces
}{%
  \end{tcolorbox}%
}

\NewDocumentEnvironment{property}{o}{%
  \refstepcounter{ltproperty}%
  \LT@FloatLocal
  \begin{tcolorbox}[ltbox, colframe=qiancheng, colback=qiancheng!7!white]
    \noindent{\bfseries\color{qiancheng}\propertyname~\theltproperty
      \IfValueT{#1}{\ （#1）}}\par\nobreak\smallskip
    \ignorespaces
}{%
  \end{tcolorbox}%
}

% --- 例题：内容盒，注（题源）直接进标题行 -----------------------------------
\NewDocumentEnvironment{example}{o}{%
  \refstepcounter{ltexample}%
  \LT@FloatLocal
  \begin{tcolorbox}[ltcontent, colframe=example, colback=example!5!white]
    \noindent{\bfseries\color{example}\examplename~\theltexample
      \IfValueT{#1}{\ （#1）}}\par\nobreak\smallskip
    \ignorespaces
}{%
  \end{tcolorbox}%
}

% --- 解：不编号，跟在例题之后，安静的中性色 ---------------------------------
\NewDocumentEnvironment{solution}{o}{%
  \par\smallskip
  \noindent{\bfseries\color{InkMuted}\solutionname\IfValueT{#1}{\ （#1）}}\,%
  \ignorespaces
}{%
  \par\smallskip
}

% ============================================================================
%  美术标题栏模式（类选项 boxart）：复刻原版孤星笔记模板的标题盒
%  左右两张预渲染矢量图（assets/pdf/<色>_1/2.pdf，含 Rhine logo、竖排中文
%  环境名、英文大字与三色条，全部为路径轮廓，无字体依赖）+ 中间色块标签。
%  盒体回到原版形制：主题色 2.5pt 粗框 + qianbai 底；盒内正文字体不改
%  （书籍正文不切楷体）。工程需有 assets/tikz 源并已渲染出 assets/pdf；
%  换主题后重跑 python tools/lonetrail.py 的渲染即同步换色。
% ============================================================================
\iflt@boxart
\newlength{\LT@artwd}
% 带名称：左图 + 名称色块 + 右图
% 注意不要写 \XeTeXinterchartokenstate=0（原笔记版有）：它会关掉 xeCJK 的
% 字符路由，色块里的中文会回落到西文字体变豆腐块。
\newcommand{\LT@ArtBar}[3]{%
  \sbox0{\fontsize{14}{22}\selectfont\heiti #2}%
  \settowidth{\LT@artwd}{\usebox0}%
  \begin{tabular}[c]{@{}c@{}c@{}c@{}}
    \raisebox{-\height+12pt}{\includegraphics[height=24pt]{assets/pdf/#1.pdf}} &
    \raisebox{-\height+12pt}{%
      \colorbox{heise}{%
        \parbox[c][18pt][c]{\LT@artwd}{%
          \centering\fontsize{14}{22}\selectfont\heiti
          \color{mibai} #2\vspace{-4pt}}}} &
    \raisebox{-\height+12pt}{\includegraphics[height=24pt]{assets/pdf/#3.pdf}}
  \end{tabular}%
}
% 不带名称：左右图直接相接（竖排中文环境名在左图里）
\newcommand{\LT@ArtBarPlain}[2]{%
  \begin{tabular}[c]{@{}c@{}c@{}}
    \raisebox{-\height+12pt}{\includegraphics[height=24pt]{assets/pdf/#1.pdf}} &
    \raisebox{-\height+12pt}{\includegraphics[height=24pt]{assets/pdf/#2.pdf}}
  \end{tabular}%
}
\tcbset{
  ltboxart/.style={
    enhanced,
    sharp corners,
    breakable,
    colback=qianbai,
    boxrule=2.5pt,
    arc=0pt,
    fontupper=\normalfont,
    halign upper=justify,
    left=10pt,
    right=10pt,
    top=6pt,
    bottom=6pt,
    before skip=14pt,
    after skip=14pt,
    attach boxed title to top left={yshift=-2.5pt, xshift=0mm},
    boxed title style={
      boxrule=0pt,
      boxsep=0pt,
      top=-2.1pt,
      bottom=0pt,
      left=0pt,
      right=0pt,
      colback=qianbai,
      baseline=\ht\strutbox
    },
  }
}
% 工厂：#1=环境名 #2=主题色名（计数器 lt<env>，名宏 <env>name 已在上文定义）
\newcommand{\LT@ArtBox}[2]{%
  \RenewDocumentEnvironment{#1}{o}{%
    \refstepcounter{lt#1}%
    \LT@FloatLocal
    \IfValueTF{##1}{%
      \begin{tcolorbox}[ltboxart, colframe=#2,
        title={\LT@ArtBar{#2_1}{##1}{#2_2}}]%
    }{%
      \begin{tcolorbox}[ltboxart, colframe=#2,
        title={\LT@ArtBarPlain{#2_1}{#2_2}}]%
    }%
    \ignorespaces
  }{\end{tcolorbox}}%
}
\LT@ArtBox{definition}{definition}
\LT@ArtBox{theorem}{theorem}
\LT@ArtBox{proposition}{proposition}
\LT@ArtBox{corollary}{corollary}
\LT@ArtBox{lemma}{lemma}
\LT@ArtBox{example}{example}
\fi

% --- 习题 -------------------------------------------------------------------
\NewDocumentEnvironment{exercise}{o}{%
  \refstepcounter{bookexercise}%
  \par\smallskip\noindent
  {\bfseries\color{main}\exercisename~\thebookexercise
    \IfValueT{#1}{\ （#1）}}\quad\ignorespaces
}{%
  \par\smallskip
}

% 节末练习块（连续编号）
\newenvironment{exercises}{%
  \par\vspace{0.6em}\noindent
  {\bfseries\color{main}\exercisename}\par\smallskip
  \ignorespaces
}{%
  \par\vspace{0.6em}%
}

% 章末练习集（标题 + 连续编号）
\newenvironment{Exercises}{%
  \markright{习题}%
  \addcontentsline{toc}{section}{习题}%
  \par\vspace{0.8em}%
  \begin{center}{\large\bfseries\color{InkBrown}\heiti 习题}\end{center}%
  \vspace{0.5em}\small
  \list{\bfseries\thebookexercise.}{%
    \usecounter{bookexercise}%
    \renewcommand{\theHbookexercise}{Exercise.\thechapter.\arabic{bookexercise}}%
  }%
}{%
  \endlist\normalsize\par\vspace{0.8em}%
}

% --- 证明 -------------------------------------------------------------------
\NewDocumentEnvironment{proof}{o}{%
  \par\smallskip\noindent
  {\bfseries\color{InkMuted}\proofname\IfValueT{#1}{\ （#1）}}\;
  \color{InkJoint}\ignorespaces
}{%
  \hfill{\color{InkMuted}$\square$}\par\smallskip
}

% ============================================================================
%  注：短 → 边栏；长 → 正文静音块（与原书阈值一致，行为不变）
% ============================================================================
\ExplSyntaxOn
\bool_new:N \l_lt_remark_long_bool
\cs_new_protected:Npn \lt_remark_classify:n #1
  {
    \bool_set_false:N \l_lt_remark_long_bool
    \int_compare:nNnT { \tl_count:n {#1} } > { 64 }
      { \bool_set_true:N \l_lt_remark_long_bool }
    \bool_if:NF \l_lt_remark_long_bool
      {
        \regex_match:nnTF
          {
            \\begin\s*\{
              ( equation|align|alignat|flalign|gather|multline|itemize|enumerate|
                compactitem|compactenum|tikzpicture|multicols|choices|figure|table|
                tabular|tabularx )\s*\*?\s*\}
            | \\\[
            | \\par\b
          }
          { \tl_to_str:n {#1} }
          { \bool_set_true:N \l_lt_remark_long_bool }
          { }
      }
  }
\cs_new_protected:Npn \lt_remark_margin:nn #1 #2
  {
    \refstepcounter{ltremark}%
    \sidenote
      {
        \textbf{\remarkname~\theltremark\tl_if_novalue:nF {#1} {\,(#1)}}
        \par #2
      }
  }
\cs_new_protected:Npn \lt_remark_main:nn #1 #2
  {
    \refstepcounter{ltremark}%
    \begin{tcolorbox}[ltremarkbox]
      \LT@FloatLocal
      {\footnotesize\bfseries\color{InkMuted}\remarkname~\theltremark
        \tl_if_novalue:nF {#1} {\,(#1)}}\par\nobreak
      {\footnotesize\color{InkMuted}\raggedright #2\par}
    \end{tcolorbox}%
  }
\cs_new_protected:Npn \lt_remark_auto:nn #1 #2
  { \lt_remark_classify:n {#2} \bool_if:NTF \l_lt_remark_long_bool
      { \lt_remark_main:nn {#1}{#2} } { \lt_remark_margin:nn {#1}{#2} } }
\cs_set_eq:NN \LTRemarkAuto \lt_remark_auto:nn
\ExplSyntaxOff

\NewDocumentEnvironment{remark}{o +b}{\LTRemarkAuto{#1}{#2}}{}
% 显式强制正文形式（长注 / 人工指定）
\NewDocumentEnvironment{remark*}{o +b}{%
  \refstepcounter{ltremark}%
  \begin{tcolorbox}[ltremarkbox]
    \LT@FloatLocal
    {\footnotesize\bfseries\color{InkMuted}\remarkname~\theltremark
      \IfValueT{#1}{\ （#1）}}\par\nobreak
    {\footnotesize\color{InkMuted}\raggedright #2\par}
  \end{tcolorbox}%
}{}

% ============================================================================
%  盒子内的浮动与边注本地化
%  tcolorbox 内不允许浮动体（"Not in outer par mode"），figure/table 就地锁 [H]；
%  \marginpar 在盒内会丢，就地降级为行内括注。
% ============================================================================
\newcommand{\LT@FloatLocal}{%
  \let\LT@saved@figure\figure
  \let\LT@saved@table\table
  \def\figure{\@ifnextchar[{\LT@fig@opt}{\LT@saved@figure[H]}}%
  \def\LT@fig@opt[##1]{\LT@saved@figure[H]}%
  \def\table{\@ifnextchar[{\LT@tbl@opt}{\LT@saved@table[H]}}%
  \def\LT@tbl@opt[##1]{\LT@saved@table[H]}%
  \let\LT@saved@sidenote\sidenote
  \def\sidenote{\@ifnextchar[{\LT@sn@opt}{\LT@sn@plain}}%
  \def\LT@sn@opt[##1]##2{\LT@sn@body{##2}}%
  \def\LT@sn@plain##1{\LT@sn@body{##1}}%
}
\newcommand{\LT@sn@body}[1]{{\footnotesize\color{InkMuted}（#1）}}

% ============================================================================
%  章首阅读提示 / 延伸讨论（正文内容盒）
% ============================================================================
\newtcolorbox{wenxintishi}[1][阅读提示]{%
  ltbox, colframe=InkMuted, colback=InkWash!45!white, leftrule=4pt,
  top=7pt, bottom=7pt, before skip=10pt, after skip=11pt,
  title={{\bfseries\color{InkBrown}\heiti #1}},
  coltitle=InkBrown, colbacktitle=InkWash!45!white,
  fonttitle=\bfseries,
  attach boxed title to top left={yshift=-\tcboxedtitleheight/2,
    xshift=\tcboxedtitlewidth/4},
  boxed title style={boxrule=0pt, colframe=InkWash!45!white,
    colback=InkWash!45!white}
}

\NewDocumentEnvironment{analysispath}{O{延伸讨论}}{%
  \par\vspace{0.7em}%
  \noindent{\sffamily\bfseries\small\color{main} #1}\par\nobreak
  \vspace{0.35em}%
  \ignorespaces
}{%
  \par\vspace{0.5em}%
}

% ============================================================================
%  选择题选项统一环境
% ============================================================================
\newlist{choices}{enumerate}{1}
\setlist[choices]{
  label = \textbf{\Alph*.},
  labelsep = 0.5em,
  leftmargin = *,
  itemsep = 0.35em,
  topsep = 0.4em,
  parsep = 0pt,
  partopsep = 0pt
}

% ============================================================================
%  习题解答收集：exsol 收体，附录按章统一输出
% ============================================================================
\newif\iflt@noanswer
\AtBeginDocument{%
  \@ifundefined{ELEGANT@result}{}{%
    \ifdefstring{\ELEGANT@result}{noanswer}{\lt@noanswertrue}{}}%
}
\def\lt@sol@chlist{}
% 追加一条解答到 per-chapter 存储宏。
% 正文 #3 必须经 toks 寄存器原样存入 —— 若直接 \xdef 会把解答体里的
% \begin{enumerate}[...]、\fpeval 等可展开物当场展开炸掉。
\newtoks\lt@sol@existing
\newtoks\lt@sol@prefixtoks
\newtoks\lt@sol@bodytoks
\newcommand{\lt@sol@append}[3]{%
  \@ifundefined{lt@sol@ch@#1}{%
    \expandafter\gdef\csname lt@sol@ch@#1\endcsname{}%
    \ifx\lt@sol@chlist\empty
      \xdef\lt@sol@chlist{#1}%
    \else
      \xdef\lt@sol@chlist{\lt@sol@chlist,#1}%
    \fi
  }{}%
  \begingroup
    \lt@sol@existing\expandafter\expandafter\expandafter{%
      \csname lt@sol@ch@#1\endcsname}%
    \lt@sol@prefixtoks\expandafter{#2}%
    \lt@sol@bodytoks{#3}%
    \expandafter\xdef\csname lt@sol@ch@#1\endcsname{%
      \the\lt@sol@existing
      \the\lt@sol@prefixtoks
      \the\lt@sol@bodytoks
    }%
  \endgroup
}
\newcommand{\lt@sol@emit}[1]{%
  \@ifundefined{lt@sol@ch@#1}{}{\csname lt@sol@ch@#1\endcsname}%
}
% 章标题捕获（解答按「第N章 章名」分节）
\let\lt@orig@makechapterhead\@makechapterhead
\def\@makechapterhead#1{%
  \if@mainmatter
    \expandafter\gdef\csname lt@sol@chtitle@\the\c@chapter\endcsname{#1}%
  \fi
  \lt@orig@makechapterhead{#1}%
}

\NewDocumentEnvironment{exsol}{+b}{%
  \iflt@noanswer\else
    \refstepcounter{exsolcount}%
    \xdef\lt@sol@prefix{%
      \noexpand\par\noexpand\noindent
      \noexpand\textbf{\noexpand\color{main}\the\c@chapter.\the\c@exsolcount.}%
      \noexpand\quad
      \noexpand\textbf{\noexpand\color{main}\noexpand\solutionname}%
      \noexpand\quad
    }%
    \lt@sol@append{\thechapter}{\lt@sol@prefix}{#1\par}%
  \fi
}{}

% 对外接口
\newcommand{\printexercisesolutions}[1]{\lt@sol@emit{#1}}
\newcommand{\printallexercisesolutions}{%
  \ifx\lt@sol@chlist\empty
  \else
    \chapter{习题解答}%
    \@for\lt@sol@curch:=\lt@sol@chlist\do{%
      \expandafter\lt@sol@emitone\expandafter{\lt@sol@curch}}%
  \fi
}
\newcommand{\lt@sol@emitone}[1]{%
  \@ifundefined{lt@sol@ch@#1}{}{%
    \section*{第#1章 \csname lt@sol@chtitle@#1\endcsname}%
    \addcontentsline{toc}{section}{第#1章 \csname lt@sol@chtitle@#1\endcsname}%
    \csname lt@sol@ch@#1\endcsname}%
}
% 回到类文件的 catcode 上下文（.cls 内 @ 为 letter）
